\documentclass{article}
\usepackage{graphicx} % Required for inserting images
\usepackage{mathtools}
\usepackage{mathrsfs}
\usepackage{amsmath}
\usepackage{amssymb}
\usepackage{bm}
\usepackage{booktabs}
\usepackage[utf8]{inputenc}
\usepackage[T1]{fontenc}
\usepackage[ruled, vlined, linesnumbered]{algorithm2e}
\usepackage[margin=2cm]{geometry}
 
\SetKwInput{KwInput}{Input}
\SetKwInput{KwOutput}{Output}
\SetKwComment{Comment}{$\triangleright$\ }{}
 
\title{Regime Detection using Sliced Wasserstein k-means clustering with UnifOrtho}
\author{Oscar Peyron \\ \small \textit{ADIA Lab - University College London}}
\date{February 2026}

\begin{document}
\maketitle
\section{Abstract}  
This thesis investigates the topic of regime Detection in high dimensional time-series using an interpretable and efficient machine Learning algorithm. 
In particular, we provide an enhanced algorithm to cluster the multi-dimensional returns distributions building upon the Wasserstein clustering algorithm \cite{HorvathEtAl2024} and its high-dimensional generalization called the Sliced Wasserstein clustering algorithm proposed by \cite{LuanHamp2025}. One of the key characteristics of those clustering algorithms  is their ability
to infer which regime the market is in without pre-specified assumptions on the returns distribution.
Specifically, they use the Wasserstein or sliced-Wasserstein (in higher dimensions) distance as a metric to compare different returns distributions. Our focus will be on improving the Sliced Wasserstein clustering algorithm in high-dimensional settings.
This sliced-Wasserstein method has been designed to project high-dimensional distributions into one-dimensional distributions and then compute the Wasserstein distance in one dimension, which is computationally efficient.
However, in order to get a good approximation of the sliced Wasserstein distance, one needs to sample a large number of random projection vectors to approximate the latter distance. 
This is exactly, the problem we are trying to solve in this thesis. Instead of using a fixed set of well-chosen random projection vectors to compute the Sliced-Wasserstein distance as in Luamp et. al \cite{LuanHamp2025}, we propose to use the UnifOrtho method designed by Rowland et. al \cite{RowlandEtAl2019} to sample orthogonal projections, which leads to a significant reduction in variance and more stable clustering results in high-dimensions.
Another, key contribution of this research is to use the clustering results from the improved Sliced-Wasserstein distance, to derive implied probabilities of regime swtiches in the market using simple softmax functions and bayesian priors. 
Finally, we design multiple trading strategies based on the above regime detection method and compare their performance using statistical tests such as the Ledoit-Wolf test on Sharpe ratios.

\section{Literature Review}
Regime detection in financial markets is a heavily studied topic in literature, as it is of great interest for practitioners. 
Notably, knowing regimes helps actively investors to adapt their trading strategies and risk management to the current market conditions. 
One of the first notable papers by Hamilton \cite{Hamilton1989} introduced a parametric approach to regime detection using hidden Markov models (HMM). 
This method is until now, one of the most widely used methods for regime detection in financial markets.
Rydén et al. \cite{Rydén1998} showed that the HMMs approach captures heavy-tails and unconditional non-normality in the returns distriubtion, in addition to being a highly intuitive method for regime detection.
More recently, Ang et al. \cite{AngTimmermann2012} showed that regime in financial markets empirically switches abruptly and this changed behavior persists for many periods. 
However, HMMs struggle as the framework assumes that regimes are a mixture of Gaussian distributions, which is not always the case in financial markets.
Moreover, breaking points in the markets forces the model to adapts its parameters for the return distribution, which can be difficult to do in practice.
Traditional regime detection approach try to identify these global beaking-points for adapting parameters such as in Bai et al.\cite{BaiPerron2003} and with general-to-specific modeling approaches in Hendry et.al \cite{Hendry2014}.
The Wasserstein -k-means approach as proposed by Horvath et al. \cite{HorvathEtAl2024} is instead classifying regime by analysing the behaviour of the \textit{distribution} of the time-series of returns through the study period.
This enables to be model-free for the returns distribution and therefore skip parameters adaptation issues.
The Sliced Wasserstein k-means clustering algorithm as proposed by Luan et al. \cite{LuanHamp2025} generalizes Horvath's approach to high-dimensional time-series data and infers regimes changes using returns distribution in $\mathcal{P}_p(\mathbb{R}^d)$, the space of probability measures on $\mathbb{R}^d$ with finite $p$-th moment.






\section{Theoretical Background}
This chapter will focus on the theoretical framework to understand the Wasserstein distance and its extension to the Sliced Wasserstein distance.
It will also provide the mathematics behind the implied probability derived from the Wasserstein distance to produce predictive signals for trading strategies. 
Finally, we will provide some background on the statistical tests we used to compare the performance of the different trading strategies (i.e Ledoit-Wolf test on Sharpe ratios). 

\subsection{Theoretical Framework for Wasserstein Distance and Wasserstein Barycenters}
In this chapter we consider $\mathcal{X} = \mathbb{R}^d$ with $d$ the dimension of the returns vector. 
\begin{paragraph}{Definition: Data Streams}
A set $\mathcal{S}$ of streams of data over $\mathcal{X}$ is defined as: $$S(\mathcal{X}) = \left\{(X_1, X_2, \dots, X_N) | X_i \in \mathcal{X}, N \in \mathbb{N} \right\}.$$ 
For instance in this research we consider the following: 
\begin{equation}
    S = (S_1, S_2, \dots, S_N) \in \mathcal{S}(\mathbb{R}^d)
\end{equation}
where $S_i$ might be the the price vector of the different assets at time $i$.
\cite{HorvathEtAl2024}
\end{paragraph}

\begin{paragraph}{Definition: Lifting Transformation}
A lifting transformation $l$ is a map that maps the set of streams $\mathcal{S}(\mathcal{X})$ into the set of streams of streams: 
\begin{equation}
l: \left\{
\begin{aligned}
\mathcal{S}(\mathcal{X}) \rightarrow \mathcal{S}(\mathcal{S}(\mathcal{X})) \\ 
S \mapsto l(S) = (l_1, l_2, \dots, l_M)
\end{aligned}
\right.
\end{equation}
where $M>1$.  \cite{HorvathEtAl2024}
\end{paragraph}

The example from \cite{BonnierKidger2019} is the sliding window transformation: 
\begin{equation}
    l^m(\mathbf{x}) = (x_{1+h_2 (m-1)}, x_{2+h_2 (m-1)}, \dots, x_{h_1+h_2 (m-1)}) \text{ for } m = 1, 2, \dots, M
\end{equation}
with $M := \left\lfloor\frac{N-h_1-h_2}{h_2}\right\rfloor$ and $h_1, h_2 \in \mathbb{N}$ the window size and the step size respectively.

\paragraph{Definition: Discrete uniform distribution}
Given a stream generated by the lifting process $l(S) = (l^m(\mathbf{x})_i, : i = 1, 2, \dots, h_1)$, we can define the discrete uniform distribution $\mu^S$ as follows:
\begin{equation}
\mu_m = \frac{1}{h_1} \sum_{i=1}^{h_1} \delta_{l^m(\mathbf{x})_i}
\end{equation}
where we have that $\delta_x$ is the Dirac delta at point $x$. \cite{LuanHamp2025}
This definition allows us to introduce the set of empirical distributions $\mathcal{K} = \{\mu_j\}_{1 \le j \le M}$ where $\mu_j \in \mathcal{P}_p(\mathbb{R}^d)$ with $\mathcal{P}_p(\mathbb{R}^d)$ the set of probability measures on $\mathbb{R}^d$ with finite $p$-th moment. \cite{LuanHamp2025}. 


\paragraph{Definition: Wasserstein distance $\mathcal{W}_p$}
Let $\mu$ and $\nu$ be two probability measures on $\mathbb{R}^d$ with finite $p$-th moment defined as the following infimum over the joint distributions $(X,Y)$ of $d$-dimensional random vectors $X$ and $Y$. The Wasserstein distance of order $p$ between $\mu$ and $\nu$ is defined as:
\begin{equation}
\mathcal{W}_p(\mu, \nu) =  \inf_{(X,Y) \sim \Gamma(\mu, \nu)} \left( \mathbb{E}\left[ \|X - Y\|^p \right] \right)^{1/p} = \left( \inf_{\pi \in \Gamma(\mu, \nu)} \int_{\mathbb{R}^d \times \mathbb{R}^d} \|x - y\|^p d\pi(x,y) \right)^{1/p}
\end{equation}
where $\Gamma(\mu, \nu) \subseteq \mathcal{P}_p(\mathbb{R}^d \times \mathbb{R}^d)$ is the set of joint distributions over $\mathbb{R}^d \times \mathbb{R}^d$ with marginals $\mu$ and $\nu$. An element $\pi^*\in \Gamma(\mu, \nu)$ that achieves the infimum is called an optimal transport plan between $\mu$ and $\nu$. \cite{RowlandEtAl2019}


\paragraph{Definition: Wasserstein barycentre $\bar{\mu}$}
Given a family probability measures $\mathcal{K}$ in $\mathbb{R}^d$ with finite $p$-th moment, the Wasserstein barycentre $\bar{\mu}$ is defined as the probability measure that minimizes the sum of squared Wasserstein distances to the given measures:
\begin{equation}
\bar{\mu} = \arg\min_{\nu \in \mathcal{P}_p(\mathbb{R}^d)} \sum_{\mu_m \in \mathcal{K}} \mathcal{W}_p(\nu, \mu_m)
\end{equation}
where $\mathcal{P}_p(\mathbb{R}^d)$ is the set of probability measures on $\mathbb{R}^d$ with finite $p$-th moment. \cite{LuanHamp2025}

\paragraph{Projected Empirical Measure}
Let $\theta \in \mathbb{S}^{d-1}$  with $\mathbb{S}^{d-1}$ the unit sphere in $d$ dimensions.  
The projected empirical measure $\mu'(\theta) \in \mathcal{P}_{p}(\mathbb{R})$ is defined as:
\begin{equation}
\mu'(\theta) = \frac{1}{N} \sum_{i} \delta_{\langle x_i, \theta \rangle}
\end{equation}
with $x_i$ the atoms of the original empirical measure $\mu \in \mathcal{P}_p(\mathbb{R}^d)$ and $\langle \cdot, \cdot \rangle$ the standard inner product in $\mathbb{R}^d$. \cite{LuanHamp2025}


\paragraph{Definition: The sliced Wasserstein distance $\mathcal{\bar{W}}_p$}
Given two distributions $\mu, \nu \in \mathcal{P}_p(\mathbb{R}^d)$ with $d>1$, the sliced Wasserstein distance is defined as follows: 
\begin{equation}
\mathcal{\widetilde{W}}_p(\mu, \nu) = \int_{\mathbb{S}^{d-1}} \mathcal{W}_p(\mu'(\theta), \nu'(\theta)) d\theta
\end{equation}
with $\mathcal{W}_p(\mu'(\theta), \nu'(\theta))$ the Wasserstein distance in d=1 between the projected empirical measures $\mu'(\theta)$ and $\nu'(\theta)$ on the direction $\theta$. \cite{LuanHamp2025}

In practice one can approximate the sliced Wasserstein distance using Monte Carlo methods by sampling $L$ random directions $\{\theta^1, \dots, \theta^L\}$ uniformly on the unit sphere $\mathbb{S}^{d-1}$ and computing the average Wasserstein distance over these projections:
\begin{equation}
\mathcal{\widetilde{W}}_p(\mu, \nu) \approx \frac{1}{L} \sum_{l=1}^{L} \mathcal{W}_p(\mu'(\theta^l), \nu'(\theta^l))
\end{equation}

More precisely we have, $\theta^1, \dots, \theta^L \sim^{i.i.d.} \text{Unif}(\mathbb{S}^{d-1})$.  \cite{RowlandEtAl2019}


\paragraph{Definition: The sliced Wasserstein Barycenter}
The sliced Wasserstein barycenter $\bar{\mu}$ of a family of probability measures $\mathcal{K} = \{\mu_j\}_{1 \le j \le M}$ is defined as:
\begin{equation}
\bar{\mu}^{\widetilde{\mathcal{W}}_p} = \arg\min_{\nu \in \mathcal{P}_p(\mathbb{R}^d)} \sum_{\mu_m \in \mathcal{K}} \widetilde{\mathcal{W}}_p^p(\nu, \mu_m)
\end{equation}
\cite{LuanHamp2025}


%%% explain more in detail was is the Wasserstein distance and the basics 


\subsection{Theoretical Framework for Sliced Wasserstein (SW) with UnifOrtho Projections}
This section will provide insights on the theoretical background behind the Uniortho approach to compute sliced Wasserstein distance in the case of the Sliced-Wasserstein k-means clustering algorithm\cite{LuanHamp2025}
Notably, we will see that this method solves the low generalizable nature of using pre-defined projection vectors ${\theta^l}$  for the Sliced Wasserstein distance as proposed by Luan et al \cite{LuanHamp2025} and Horvath et al \cite{HorvathEtAl2024}. Instead the Unifortho approach provides a more robust and generalizable clustering algorithm in higher-dimensions. 

As we have seen in the previous section, the Sliced Wasserstein distance is approximated using $L$ sampling of random direction in the unit sphere $\mathbb{S}^{d-1}$.
The UnifOrtho improvement proposed in \cite{RowlandEtAl2019} and \cite{PetrovicBardenetDesolneux2025} is to sample an entire $N$ size vector from the $\text{UnifOrt}(\mathbb{S}^{d-1}, N)$ distribution which is the distribution of $N$-tuples of orthogonal vectors uniformly distributed on the unit sphere.


It provides the following advantages compared to standard Monte Carlo methods. 
\begin{itemize}
    \item \textbf{Variance reduction via mutual orthogonality} : The mutual orthogonality of the sampled directions in the UnifOrtho method leads to a significant reduction in the variance of the sliced Wasserstein distance estimator compared to standard Monte Carlo methods that sample independent random directions. This is because the orthogonal directions provide more diverse and informative projections of the distributions, leading to a more accurate estimation of the sliced Wasserstein distance.
    \item \textbf{Superior performance in high dimensions} : The UnifOrtho is therefore perfectly adapted to working with multi-dimensional time series for the SW k-means clustering algorithm.
    \item \textbf{Clustering Stability}: The reduced variance of the sliced Wasserstein distance estimator leads to more stable clustering results, as the algorithm is less sensitive to the randomness in the projection directions. This stability is particularly beneficial in high-dimensional settings where the curse of dimensionality can exacerbate the variability of the estimator.
\end{itemize}


\paragraph{Definition: Orthogonal group $O(d)$}
The orthogonal group $O(d)$ is defined as :
\begin{equation}
    O(d) = \{ U \in \mathbb{R}^{d \times d} \mid U^T U = I_d \}
\end{equation}

\paragraph{Definition: Haar measure on $O(d)$}
The Haar measure on $O(d)$ is the unique probability measure that satisfies the following properties:
\begin{enumerate}
    \item $\mu(QA) = \mu(A)$ for all $Q \in O(d)$ and Borel sets $A \subseteq O(d)$ (left-invariance)
    \item $\mu(AQ) = \mu(A)$ for all $Q \in O(d)$ and Borel sets $A \subseteq O(d)$ (right-invariance)
    \item $\mu(O(d)) = 1$ (normalization)
\end{enumerate}


\paragraph{Definition: $\text{UnifOrtho}(\mathbb{S}^{d-1}, N)$}
Let $N \le d$ be a positive integer. The probability distribution $\text{UnifOrtho}(\mathbb{S}^{d-1}, N) \in \mathcal{P}((\mathbb{S}^{d-1})^N)$ is defined as the joint distribution of $N$ rows of a random matrix drawn from the Haar distribution on the orthogonal group $O(d)$. In other words, if $U$ is a random matrix drawn from the Haar distribution on $O(d)$, then the distribution of any $N$ rows of $U$ is $\text{UnifOrtho}(\mathbb{S}^{d-1},N)$. \cite{RowlandEtAl2019}

\paragraph{Numerical Recipe for constructing matrices with Haar distribution on $O(d)$}
\begin{enumerate}
    \item Take an $N \times N$ complex matrix $Z$ whose entries are complex standard 
          normal random variables.
    
    \item Feed $Z$ into \textit{any} QR decomposition routine. Let $(Q, R)$, where 
          $Z = QR$, be the output.
    
    \item Create the following diagonal matrix
    \begin{equation}
        \Lambda = 
        \begin{pmatrix}
            \dfrac{r_{11}}{|r_{11}|} & & \\[6pt]
            & \ddots & \\[6pt]
            & & \dfrac{r_{NN}}{|r_{NN}|}
        \end{pmatrix},
        \label{eq:5.12}
    \end{equation}
    where the $r_{jj}$s are the diagonal elements of $R$.
    
    \item The diagonal elements of $R' = \Lambda^{-1}R$ are \textit{always} real and 
          strictly positive, therefore the matrix $Q' = Q\Lambda$ is distributed 
          with Haar measure.
\end{enumerate}\cite{Mezzadri2007}




The following pseudocode is based on the work by \cite{RowlandEtAl2019,Mezzadri2007} on UnifOrtho methods and Haar Measure. The general approach is inspired from \cite{LuanHamp2025} for the Sliced Wasserstein k-means clustering for regime detection. 

\begin{algorithm}[H]
\caption{sWk-means clustering with UnifOrtho Projections}\label{alg:swkmeans}

\KwInput{Data stream $S \in \mathcal{S}(\mathbb{R}^d)$,\; number of clusters $K$,\; number of projections $L$,\; tolerance $\epsilon$ for convergence}
\KwOutput{$K$ cluster assignments and corresponding 1d barycenters}

\BlankLine

\Comment{Preprocessing}
Calculate $l(r^S)$ given $S \in \mathcal{S}(\mathbb{R}^d)$ and define family of empirical distributions $\mathcal{K} = \{\mu_j\}_{1 \le j \le M}$\;

\BlankLine
\Comment{Step 1: UnifOrtho Projection Generation}
$k \leftarrow \lceil L/d \rceil$ \Comment*[r]{number of orthogonal bases needed}
$\Theta \leftarrow \emptyset$\;

\For{$i = 1$ \KwTo $k$}{
    Sample a random matrix $Z_i \in \mathbb{R}^{d \times d}$ with i.i.d.\ $\mathcal{N}(0,1)$ entries\;
    Compute the QR decomposition: $Z_i = Q_i R_i$\;
    Let $\Lambda_i$ be a diagonal matrix where $\Lambda_{i, jj} = \text{sgn}(R_{i, jj})$\;
    $U_i \leftarrow Q_i \Lambda_i$ \Comment*[r]{$U_i$ is Haar-distributed on $O(d)$}
    Add the $d$ column vectors of $U_i$ to $\Theta$\;
}
$\{\theta^1, \dots, \theta^L\} \leftarrow \Theta[1:L]$ \Comment*[r]{select exactly $L$ directions}

\BlankLine
\Comment{Step 2: Projection \& K-Means Iteration}
Calculate projected distributions $\{\mu_j'(\theta^l)\}_{1 \le j \le M}$ for all $l=1 \dots L$\;
Initialize cluster assignments $C_1, \dots, C_K$\;

\While{convergence criterion $> \epsilon$}{
    \For{$k = 1$ \KwTo $K$}{
        Compute the 1d Wasserstein barycenter $\bar{\mu}_k(\theta^l)$ for each projection $l$ using empirical distributions assigned to $C_k$\;
    }
    \For{$j = 1$ \KwTo $M$}{
        Assign $\mu_j$ to the cluster $C_{k^*}$ where $k^* = \arg\min_{k} \sum_{l=1}^L W_p^p\!\left(\mu_j'(\theta^l),\, \bar{\mu}_k(\theta^l)\right)$\;
    }
}

\Return{$C_1, \dots, C_K$ and barycenters $\{\bar{\mu}_k(\theta^l)\}$}
\end{algorithm}

\paragraph{Time Complexity}
The Time Complexity of the UnifOrtho projection generation step is $\mathcal{O}(d^3)$ due to the QR decomposition of $k$ matrices of size $d \times d$. 
The projection and k-means iteration step has a complexity of $\mathcal{O}(M K L)$ for the assignment step and $\mathcal{O}(K L M_k \log M_k)$ for computing the 1d Wasserstein barycenters, where $M_k$ is the number of distributions assigned to cluster $C_k$. 


\textbf{One improvement that could be done is to use the Hadamard product, for reducing the time complexity to $\mathcal{O}(d^2 \log d)$ for generating the UnifOrtho projections by leveraging structured random matrices } \cite{ChoromanskiEtAl2019}.

\subsection{Theoretical Framework for Implied Probability Regime Detection Strategy}

\begin{paragraph}{Theoretical Background}
We use implied probability measure to compute wether on thinks the market is going to become bear or bull. 
\paragraph{Likelihood Function}
It uses SW-distance of an empirical distribution $x_m$ to every centroid and then convert the likelihood via a softmax over the negative distances. 
We introduce a temperature parameter $\tau$. 
$$
\mathcal{L}(\text{regime } k \mid x_m) = \frac{e^{-d(x_m, c_k) / \tau}}{\sum_{j} e^{-d(x_m, c_j) / \tau}}
$$
with 
$$ x_m = 1/N \sum_{i=1}^{N} \delta_{l^m(\mathbf{x})_i}$$
the empirical distribution of the current window in $\mathcal{P}_p(\mathbb{R}^d)$ with a lifting transformation. And after having run the Sliced-Wasserstein k-means clustering algorithm, on the target window of the data, we have for $k = 1, \dots, K$: 
\begin{equation}
    c_k^* = \arg\min_{\nu \in \mathcal{P}_p(\mathbb{R}^d)} 
    \sum_{\mu_m \in C_k} \widetilde{W}_p^p(\nu, \mu_m), 
    \quad k = 1, \ldots, K
\end{equation}
and we have for $\nu$ and $\mu$ in $\mathcal{P}_p(\mathbb{R}^d)$:
$$
d(\nu, \mu) = \widetilde{\mathcal{W}}_p(\nu, \mu)
$$

\paragraph{Empirical Transition Matrix Markov Prior}
A market in regime i is highly likely to stay in regime i tomorrow. We model this as a first-order Markov Chain using the historical labels. 
Probability of switching from i to j is: 
$$
\pi_{i,j} = P(Z_m = j \mid Z_{m-1} = i) = \frac{\text{Count}(i \to j)}{\sum_{x=1}^{K} \text{Count}(i \to x)}
$$
If the market was $z_{current}$ our prior belief that state M is going to be at regime $k$ is : 
$$
\mathbb{P}_{prior}(k) = \pi_{z_{current},k}
$$

Then we simply use Bayes rule to derive the following :
$$
\mathbb{P}_{Bayes}(k) = \frac{\mathbb{P}_{prior}(k) \cdot \mathcal{L}(k \mid X_m)}{\sum_{j=1}^{K} \mathbb{P}_{prior}(j) \cdot \mathcal{L}(j \mid X_m)}
$$

\paragraph{Switch Probability}
Finally the switch probability does the following : 
$$
\mathbb{P}_{switch} = 1- \mathbb{P}_{Bayes}(z_{current})
$$

\paragraph{Use Gradient} 
One may also want to take into account the fact that one is moving away or getting closer to a certain regime. 

We consider a lookback window of size $\omega$, and one computes the linear trend slope of the Wasserstein distance to each centroid $k$. 
Let $t \in \{1,2, \dots, \omega \}$
We get using  OLS, $\beta_{k}$:

$$
\beta_k = \frac{\text{Covariance}(t, \, \mathcal{W}_2(X_{M-w+t}, \, C_k))}{\text{Variance}(t)}
$$

\begin{itemize}
    \item if $\beta_{k} < 0$ the distance to regime k is shrinking (we are getting close to regime k)
    \item if $\beta_{k} > 0$ the distance to regime k is growing  (we are drifting away from regime k)
\end{itemize}

We apply softmax to get a probability, with again a temperature parameter $\tau_{grad}$:
$$
P_{grad}(k) = \frac{\exp\left(-\frac{\beta_k}{\tau_{grad}}\right)}{\sum_{j=1}^{K} \exp\left(-\frac{\beta_j}{\tau_{grad}}\right)}
$$

The final switch probability using gradient is therefore: 

$$
\mathbb{P}_{switch} = 1- \left[(1-\alpha)\mathbb{P}_{Bayes}(z_{current}) + \alpha \mathbb{P}_{Grad}(z_{current}) \right]
$$

with $\alpha \in [0:1]$ a weight that can be tuned on a validation set.

\end{paragraph}

\subsection{Theoretical Framework for the Ledoit-Wolf Sharpe ratio test}

% To be added
Need to read the paper \cite{LedoitWolf2008}. 



\section{Design methodology for Monte Carlo for SW - Synthetic Data Testing}
The focus of this chapter will be on testing the algorithm on syntethic data to understand the behaviour of the algorithm and its performance in different settings.
The algorithm will be tested on 3 dimensional data followed by 4, 5 dimensional.  Performance between original and UniOrtho method will be compared and tests will span from data with multiple correlation factors 
and with diffrent target number of clusters. Finally, in terms of time-horizon, the study will focus on the equivalent of  20 years synthetic data with 7 observations a day (252 trading days a year) to be as close as possible to real world data.


Different performance metrics will be used for the regime detection : 
\begin{itemize}
    \item Total Accuracy (TA)
    \item Mean centroid-centroid distance (MCCD)
    \item Point-to-centroid distance (PCD) 
\end{itemize}

The synthetic data will be generated similarly to previous papers \cite{HorvathEtAl2024} and \cite{LuanHamp2025}
using Geometric Brownian Motion (GBM) paths $S_t/S_{t-1} = \exp(r^S_t)$ with $S_t = (S_t^1, \dots, S_t^d)$ and $r^S_t = (r^{S^{(1)}}_{t}, \dots, r^{S^{(d)}}_{t})$ 
where : 
\begin{equation}
    r^{S^{(i)}}_{t} \sim \mathcal{N}\left((\mu-\sigma^2/2)\bm{1}dt, \bm{\Sigma} dt \right)
\end{equation}.
with 
\begin{equation}
\Theta = (\mu, \sigma^2)
\end{equation}

And the covariance matrix $\bm{\Sigma}$ is defined as follows :
\begin{equation}
\Sigma_{ij} = \sigma^2 \left( \delta_{ij} + (1- \delta_{ij})\rho \right)
\end{equation}
The $\delta_{ij}$ is the Kronecker delta, $\sigma^2$ is the variance of the returns and $\rho$ is the correlation between the different assets.

Assuming, we want to cluster two different regimes (K=2), we define, as in previous papers \cite{LuanHamp2025}, Bull and Bear regimes with the following parametrization : 
\begin{itemize}
    \item $\Theta_{\text{bull}} = (0.02, 0.2)$
    \item $\Theta_{\text{bear}} = (-0.02, 0.3)$
\end{itemize}
For K = 3, similarly to \cite{LuanHamp2025} we will consider the \textit{moon-shape structure} generated using \texttt{datasets.make\_moons} function from the \texttt{sklearn} library.

This is a summary similary to the one provided in \cite{LuanHamp2025} for the different market settings on which we are going to test the sWk-means clustering algorithm with UnifOrtho projections :

\begin{table}[h]
    \centering
    \caption{Summary of synthetic dataset parameters with $d\in \left\{3,4,5\right\}$}
    \label{tab:dataset_parameters}
    \begin{tabular}{lccc}
        \toprule
        Type & Regime Bull & Regime Bear & Regime III (Bear/Moon)  \\ 
        \midrule
        Type A (K=2) & $\Theta_{\text{bull}}; \rho = +1/2$ & $\Theta_{\text{bear}}; \rho = +1/2$ & N/A \\
        Type B (K=2) & $\Theta_{\text{bull}}; \rho = +1/2$ & $\Theta_{\text{bull}}; \rho = -1/2$ & N/A \\
        Type C (K=3) & $\Theta_{\text{bull}}; \rho = +1/2$ & $\Theta_{\text{bear}}; \rho = +1/2$ & $\Theta_{\text{bear}}; \rho = -1/2$ \\
        Type D (K=3) & $\Theta_{\text{bull}}; \rho = +1/2$ & $\Theta_{\text{bear}}; \rho = +1/2$ & $\Theta_{\text{moon}}; \rho = +1/2$ \\
        \bottomrule
    \end{tabular}
\end{table}


\section{Application of sliced-SW k-means clustering with UnifOrtho projections on real world data sets}
Finally, the chapter will focus on the application of the novel algorithm on real world data sets for different types of financial markets. 
As above, the time span, wil focus on 14 to 30 years of training data (1990- 2020) and 5 years of testing data (2020- 2025). The target data sets include the following : 
\begin{itemize}
    \item FX Spot Series Data (4D)  
        \subitem EURUSD, AUDUSD, USDJPY, USDCHF, USDCAD (Major Currency Cluster) (15 min data)
    \item Equity Series Data (3D)
        \subitem S\&P 500, NASDAQ, Dow Jones (US Equity Cluster) (daily data)
        \subitem CAC 40, DAX40, FTSE 100 (European Equity Cluster) (daily data)
        \subitem Nikkei 225, Hang Seng, Kospi 200 (Asian Equity Cluster) (daily data)
    %\item  Commodity Series Data (4D and 5D)
    %    \subitem WTI Crude Oil, Brent Crude Oil, Natural Gas, Heating Oil. 
    %    \subitem Copper Aluminum, Zinc, Nickel, Lead. 
\end{itemize}

\subsection{Performance analysis of regime detection using trading strategies}
\begin{itemize}
\item Long/Short strategy based on majority regime detection signals
\item Long/Short strategy based on Implied probability signals (using the distance to the centroids as a proxy for the confidence of the regime detection)
\item Ensemble Methods: combining the above two strategies using a weighted average of the signals or an adaptative weighting scheme to generate the final trading signal.
\end{itemize}

Tested trading strategies on FX and Equity data sets. Performance is compared to a long only and short-only approach. 
Moreover, for equities, we look at the performance for different clusters (US, Europe, Asia) to see if the performance of the regime detection is better in some markets than others.
We also test, trading algorithm that use american regime detection to trade in Asian and European markets to see if the regime detection is more generalizable in some markets than others.

\subsubsection{Standard SW k-means trading strategy}
Pseudo code for this strategy is provided here : 


\begin{algorithm}[H]
\caption{\textsc{Long\_Strat\_Unifortho} --- Binary Regime Strategy}\label{alg:unifortho}
 
\KwInput{Initial capital $C_0$,\; number of simulations $N_S$,\; price matrix $S\in\mathbb{R}^{T\times d}$,\; trimming parameters $h_1, h_2$,\; window size $w$,\; number of clusters $K$,\; risk metric $m$,\; majority look-back $\ell$}
\KwOutput{Portfolio value path $V$}
 
\BlankLine
 
$R \leftarrow \text{PctReturns}(S)$ \Comment*[r]{$R_t = S_t / S_{t-1} - 1$}
$\theta \leftarrow \mathbf{1}_d / \sigma_d$ \Comment*[r]{inverse vol weighted returns}
$r^{(\text{pf})} \leftarrow R\,\theta$ \Comment*[r]{portfolio return series}
$V \leftarrow [C_0]$\;
$n \leftarrow \lfloor T / w \rfloor$\;
 
\BlankLine
\For{$i = 0$ \KwTo $n - 2$}{
    $s \leftarrow i \cdot w$, \quad $e \leftarrow (i+1)\cdot w$\;
    $W \leftarrow S[s : e]$ \Comment*[r]{current analysis window}
 
    \BlankLine
    \Comment{Step 1 --- Wasserstein-based regime detection}
    $\hat{F},\; \boldsymbol{\mu},\; \boldsymbol{\lambda} \leftarrow \textsc{MaxMCCD\_UnifOrtho}(N_S,\, W,\, K,\, L,\, \varepsilon,\, h_1,\, h_2,\, m)$\;
 
    \BlankLine
    \Comment{Step 2 --- Exponentially weighted regime classification}
    \eIf{$\ell > |\boldsymbol{\lambda}|$}{
        $\tilde{\boldsymbol{\lambda}} \leftarrow \boldsymbol{\lambda}$\;
    }{
        $\tilde{\boldsymbol{\lambda}} \leftarrow \boldsymbol{\lambda}_{[-\ell:]}$\;
    }
    \For{$k = 0, 1, \ldots, |\tilde{\boldsymbol{\lambda}}| - 1$}{
        $w_k \leftarrow \exp\!\left(-\dfrac{\ln 2}{\tau}\,(|\tilde{\boldsymbol{\lambda}}| - 1 - k)\right)$\;
    }
    $\hat{k} \leftarrow \arg\max_k \sum_{i=0}^{|\tilde{\boldsymbol{\lambda}}|-1} w_i \,\mathbf{1}\!\left[\tilde{\lambda}_i = k\right]$\;

    \BlankLine
    \Comment{Step 3 --- Generate binary signal}
    \eIf{$\hat{k} = 1$}{
        $\sigma \leftarrow +1$ \Comment*[r]{Bullish $\Rightarrow$ Long}
    }{
        $\sigma \leftarrow -1$ \Comment*[r]{Bearish $\Rightarrow$ Short}
    }
 
    \BlankLine
    \Comment{Step 4 --- Compound portfolio over next window}
    \For{each $r_t \in r^{(\mathrm{pf})}[e : e + w]$}{
        $V.\text{append}\!\Big(V[-1]\;\times\;\big(1 + \sigma\, r_t\big)\Big)$\;
    }
}
 
\Return{$V$}
\end{algorithm}
 
\vspace{1.5cm}

\subsubsection{Implied Probability SW k-means trading strategy}

\begin{algorithm}[H]
\caption{\textsc{Long\_Strat\_Implied} --- Implied-Probability Regime Strategy}\label{alg:implied}
 
\KwInput{Initial capital $C_0$,\; simulations $N_S$,\; price matrix $S\in\mathbb{R}^{T\times d}$,\; trimming $h_1, h_2$,\; window $w$,\; clusters $K$,\; metric $m$,\; signal type $\tau \in \{\texttt{continuous},\, \texttt{hysteresis},\, \texttt{conviction}\}$,\; entry threshold $\alpha_e$,\; hold threshold $\alpha_h$,\; look-back $\ell$,\; gradient flag and weight $(\texttt{g},\, \omega_g)$}
\KwOutput{Portfolio value path $V$,\; signal history $\Sigma$}
 
\BlankLine
 
$R \leftarrow \text{PctReturns}(S)$, \quad $\theta \leftarrow \mathbf{1}_d \cdot 1/\sigma_d$
%inverse vol weighted returns
\quad $r^{(\text{pf})} \leftarrow R\,\theta$\;
$V \leftarrow [C_0]$, \quad $\Sigma \leftarrow [\,]$, \quad $n \leftarrow \lfloor T/w \rfloor$\;
\BlankLine
\For{$i = 0$ \KwTo $n - 2$}{
    $s \leftarrow i\cdot w$, \quad $e \leftarrow (i+1)\cdot w$, \quad $W \leftarrow S[s:e]$\;
 
    \BlankLine
    \Comment{Step 1 --- Regime detection (same as Algorithm~\ref{alg:unifortho})}
    $\hat{F},\;\boldsymbol{\mu},\;\boldsymbol{\lambda} \leftarrow \textsc{MaxMCCD\_UnifOrtho}(N_S,\,W,\,K,\,L,\,\varepsilon,\,h_1,\,h_2,\,m)$\;
 
    \BlankLine
    \Comment{Step 2 --- Implied transition probabilities}
    $\mathbf{P},\; p_{\text{switch}},\; \mathbf{T},\; \boldsymbol{\pi} \leftarrow \textsc{ComputeImpliedProba}(\hat{F},\;\boldsymbol{\mu},\;\boldsymbol{\lambda},\;\ell,\;\texttt{g},\;\omega_g)$\;
    $\hat{k} \leftarrow \arg\max_k\;\text{count}(\boldsymbol{\lambda}_{[-h_2:]} = k)$\;
 
    \BlankLine
    \Comment{Step 3 --- Signal generation (depends on $\tau$)}
    \uIf{$\tau = \texttt{continuous}$}{
        $\sigma \leftarrow \pi_{\text{bull}} - \pi_{\text{bear}}$\;
        \If{$|\sigma| < \delta$}{$\sigma \leftarrow 0$} \Comment*[r]{dead zone}
    }
    \uElseIf{$\tau = \texttt{hysteresis}$}{
        $\sigma_{\text{prev}} \leftarrow \Sigma[-1]$ \textbf{or} $0$\;
        \uIf{$\hat{k} = \text{Bull}$}{
            \lIf{$p_{\text{switch}} \geq \alpha_e$ \textbf{and} $\sigma_{\text{prev}} \geq 0$}{$\sigma \leftarrow -1$}
            \lElseIf{$p_{\text{switch}} < \alpha_h$ \textbf{and} $\sigma_{\text{prev}} < 0$}{$\sigma \leftarrow +1$}
            \lElse{$\sigma \leftarrow 0$}
        }
        \Else(\Comment*[f]{$\hat{k} = \text{Bear}$}){
            \lIf{$p_{\text{switch}} \geq \alpha_e$ \textbf{and} $\sigma_{\text{prev}} \leq 0$}{$\sigma \leftarrow +1$}
            \lElseIf{$p_{\text{switch}} < \alpha_h$ \textbf{and} $\sigma_{\text{prev}} > 0$}{$\sigma \leftarrow -1$}
            \lElse{$\sigma \leftarrow 0$}
        }
    }
    \ElseIf{$\tau = \texttt{conviction}$}{
        $d_k \leftarrow \begin{cases} +1 & \hat{k} = \text{Bull} \\ -1 & \hat{k} = \text{Bear}\end{cases}$\;
        $c \leftarrow 1 - 1.5\, p_{\text{switch}}$ \Comment*[r]{conviction decays with switch risk}
        $\sigma \leftarrow d_k \cdot c$\;
    }
    $\Sigma.\text{append}(\sigma)$\;
 
    \BlankLine
    \Comment{Step 4 --- Compound portfolio over next window}
    \For{each $r_t \in r^{(\mathrm{pf})}[e : e+w]$}{
        $V.\text{append}\!\Big(V[-1]\;\times\;\big(1 + \sigma\, r_t\big)\Big)$\;
    }
}
\Return{$V,\;\Sigma$}
\end{algorithm}


\subsubsection{Ledoit-Wolf test on sharpe ratios}
Conducted Ledoit-Wolf test to compare the Sharpe ratios of the different trading strategies and see if the differences in performance are statistically significant.
Conducted test with a confidence level of 95\% and reported the p-values for each pairwise comparison between the trading strategies and the long-only benchmark.
Used hac method as well as the studentized circular bootstrap approach. 

\subsection{FX Data}
I have tested different trading strategies based on regime detection using the sWk-means clustering algorithm with UnifOrtho projections and the implied probability strategy derived from the SW distance to the centroids and compared it to a long-only benchmark on the
training period from 2006 to 2020 (FX data). 


\subsection{Equity Data}
I have tested different trading strategies based on regime detection using the sWk-means clustering algorithm with UnifOrtho projections and the implied probability strategy derived from the SW distance to the centroids and compared it to a long-only benchmark on the
training period from 1990 to 2020 (3483 observations for the American Equity Data).

\paragraph{American Indices Equity}
No significant statistically difference in sharpe ratios betweent the trading strategies and a long-only benchmark. 


\subsection{Tuning the hyperparameters for the trading strategies}
Tested multiple values for the hyperparameters of the trading strategies, such as the number of projection vectors L, the entry and hold thresholds, the half-life parameter for the exponentially weighted majority voting scheme, the weights for the ensemble methods and the gradient weight for the implied probability strategy with gradient.

\paragraph{Definition: Hit Ratio}
The hit ratio is defined as follows: 
\begin{equation}
    \text{Hit Ratio} = \frac{\text{Number of profitable trades}}{\text{Total number of trades}} 
\end{equation}
We use the hit ratio as a performance metric to evaluate the effectiveness of the trading strategies specifically when evaluate the different parameters for tau, or the different parameters for the half-life. A higher hit ratio indicates a higher proportion of profitable trades, which is desirable for a successful trading strategy.

\paragraph{Definition: Win/Loss Ratio}
The win/loss ratio is defined as follows:
\begin{equation}
    \text{Win/Loss Ratio} = \frac{\text{Average profit of winning trades}}{\text{Average loss of losing trades}}
\end{equation}
We will also be using this metric to evaluate the performance of the trading strategies. 


\subsubsection{Tuning the half-life parameter for the exponentially weighted majority voting scheme}
Tested multiple values for the half-life parameter $\tau$ in the exponentially weighted majority voting scheme for 
\begin{itemize}
    \item the standard long/short unifortho strategy and
    \item the implied probability strategy with hysteresis, continuous and conviction variations.
    \item the ensemble method with equally weighted signals.
\end{itemize}
The values tested are the following
\begin{equation}
    \tau \in \{1.0, 3.0, 5.0, 10.0,20.0\}
\end{equation}
The tau value is selected based on its hit ratio and win/loss ratio during the backtesting period (2006-2020).
Results: 
\begin{itemize}
    \item Long/Short standard : $\tau = 5.0$
    \item  Implied Strategies : $\tau = 5.0$ 
    \item Ensemble method : $\tau = 5.0$
\end{itemize}


\subsubsection{Tuning the weights for ensemble methods}

Tested multiple combination of weights for the ensemble methods, which combines the binary regime long/short stratey and the implied probability strategy with different variations (hysteresis, continuous, conviction).
The weights are ordered as following :
\newline [$w_{binary}$, $w_{implied\_probability(hysteresis)}$ $w_{implied\_probability(continuous)}$, $w_{implied\_probability(conviction)}$]. 
We have tested the following weights: 
\begin{itemize}
    \item \text{[0.25, 0.25, 0.25, 0.25]},
    \item \text{[0.30, 0.10, 0.30, 0.30]},
    \item \text{[0.20, 0.30, 0.20, 0.30]},
    \item  \text{[0.30, 0.20, 0.30, 0.20]},
    \item  \text{[0.10, 0.40, 0.10, 0.40]}
\end{itemize}

\subsubsection{Tuning the gradient weight for implied probability methods}
Tested multiple values for the gradient weight $\alpha$ in the implied probability strategy with gradient.
The values tested are the following 
\begin{equation}
    \alpha \in \{0.0, 0.1, 0.2, 0.3, 0.4, 0.5, 0.6, 0.7, 0.8, 0.9, 1.0\}
\end{equation}

Optimal value for $\alpha$ is determined based on the highest Sharpe ratio achieved during the backtesting period.
For the American equity data we have determined that the optimal value is $\alpha = 0.7$. As most of the variants of implied porbability strategies (continus and hysteresis), have a peak performance at $\alpha = 0.7$. 

$\alpha = 0.0$ corresponds to a pure bayesian approach. It captures current cluster membership, regime persistence and more generally where the market is. 
$\alpha = 1.0$ corresponds to a pure gradient approach. It captures the direction of the market (getting closer or further from a certain regime) and therefore captures the momentum of the market.

Hence an optimal value of $\alpha = 0.7$ suggests that  regime momentum is slightly more important than the current regime membership and regime persistence for predicting the next day return of the market.
We could also thinkg that $\mathbb{P}_{Bayers}$ and $\mathbb{P}_{Grad}$ are not redundant and capture different information about the market. 


\subsubsection{Tuning the temperature parameter $\tau$ and $\tau_{grad}$ for the implied probability strategy}
Peformed grid search for $\tau$ taken as adaptative to the switch probability as well as for $\tau_{grad}$ taken as adaptative to the gradient values.
The optimal values for $\tau$ and $\tau_{grad}$ are determined based on the highest hit ratio achieved during the backtesting period (2006-2020). 
For the American equity data, we have determined that the optimal value for the couple $(\tau, \tau_{grad})$ must be around $(0.5, 0.1)$ overall based on all the trading strategies. 

\section{Experiment 5: Interpretability Analysis of Sliced Wasserstein Regime Detection Algorithm}

We label the regimes using weekly equity data (S\&P 500, NASDAQ, Dow Jones) on the period 2006 to 2020. 
We further consider the following macro feature set:
\begin{itemize}
    \item VIX Index
    \item 10-year Treasury Yield
    \item 5-year Treasury Yield
    \item 30-year Treasury Yield
    \item Fed Funds Rate
    \item Manufacturing PMI
    \item Manufacturing Employment
    \item DXY Dollar Index
    \item Gold Spot Price
    \item Brent Crude Oil Spot Price
\end{itemize}

After having regrouped the data from the following macro-indicators from the appropriate time period, we fit thefeature set to the regime labels using the following models: 
\begin{itemize}
    \item OLS Linear Regression (probabilistic target (Probability of being in bear regime))
    \item Logistic Regression (binary target (Bull/Bear))
    \item Random Forest Classifier (binary target (Bull/Bear))
\end{itemize}
In particular, we are analysing the feature importance of the different macro-indicators during the bear regimes using the SHAP analysis toolbox. 
We observe that in particular during predicted bear regimes by the sliced-Wasserstein k-means clustering algorithm, the VIX index is a particularly important feature in the regime classification for all three models. 



\newpage
\section{Appendix}
\subsection*{The Mathematical Structure of Hadamard-Rademacher Matrices}

\textbf{Definition} (from the paper \cite{ChoromanskiEtAl2019}) The Hadamard-Rademacher chain on $O(2^L)$ is 
the Markov chain $(X_T)_{T=1}^\infty$ defined by:

\begin{equation}
    X_T = \prod_{t=1}^{T} HD_t
\end{equation}

where $(D_t)_{t=1}^\infty$ are independent random diagonal matrices with each diagonal 
element a Rademacher (i.e.\ $\mathrm{Unif}(\{+1,-1\})$) random variable, and $H$ is 
the normalized Hadamard matrix, defined as the Kronecker product:

\begin{equation}
    H = \underbrace{
        \begin{pmatrix} 
            \frac{1}{\sqrt{2}} & \frac{1}{\sqrt{2}} \\[6pt] 
            \frac{1}{\sqrt{2}} & -\frac{1}{\sqrt{2}} 
        \end{pmatrix} 
        \otimes \cdots \otimes 
        \begin{pmatrix} 
            \frac{1}{\sqrt{2}} & \frac{1}{\sqrt{2}} \\[6pt] 
            \frac{1}{\sqrt{2}} & -\frac{1}{\sqrt{2}} 
        \end{pmatrix}
    }_{L \text{ times}}
\end{equation}

These matrices $X_T$ (typically with $T \in \{1, 2, 3\}$) have been used in 
dimensionality reduction, kernel approximation, and locality-sensitive hashing.

In practice one uses $T = 3$, giving the matrix:

\begin{equation}
    M = HD_3 \cdot HD_2 \cdot HD_1
\end{equation}

where each $D_i = \mathrm{diag}(\epsilon_1^{(i)}, \ldots, \epsilon_d^{(i)})$ with 
$\epsilon_j^{(i)} \overset{\mathrm{i.i.d.}}{\sim} \mathrm{Unif}(\{+1,-1\})$.
\newpage
\bibliographystyle{plain}
\bibliography{main}

\end{document}